\section*{Глава 1. Анализ предметной области}
\addcontentsline{toc}{section}{Глава 1. Анализ предметной области}
\stepcounter{section}

\subsection{Обзор существующих платформ публикации игр}

На современном рынке платформ для публикации и распространения игр сложились две
основные категории решений. \textbf{Открытые инди-платформы} (itch.io, GameJolt)
снижают порог входа и поддерживают браузерные форматы, однако уступают крупным
площадкам по охвату аудитории. \textbf{Порталы браузерных HTML5-игр}
(Newgrounds, Яндекс Игры) специализируются на лёгком контенте, доступном без
установки, но ограничивают поддерживаемые форматы и требуют интеграции с
проприетарными SDK.

\textbf{itch.io} --- наиболее открытая из рассматриваемых платформ и наиболее
близкая по концепции к разрабатываемой системе. Поддерживает загрузку архивов любого
формата, в том числе Unity WebGL сборок для запуска непосредственно в браузере через
встроенный iframe-плеер. Разработчик сам устанавливает цену (в том числе «заплати
сколько хочешь») и долю отчислений от 0 до 30\%. Жёсткие требования к модерации
отсутствуют. Среди недостатков --- отсутствие специализированного структурированного
каталога с агрегированной статистикой, слабая ролевая модель разработчик/игрок
и обобщённость платформы без акцента на конкретные технологии игростроения.

\textbf{GameJolt} --- инди-платформа с выраженной геймификацией: система трофеев,
таблицы лидеров и достижения профиля формируют дополнительную вовлечённость аудитории.
Поддерживаются HTML5-игры. Коммерческая монетизация для разработчиков отсутствует.
Аудитория значительно меньше, чем у itch.io; инструменты аналитики ограничены.

\textbf{Newgrounds} --- один из старейших порталов браузерного контента (основан в
1995~г.), исторически ориентированный на Flash-анимации и игры. После прекращения
поддержки Flash в 2020~г. перешёл на HTML5. Обладает развитым сообществом и гибкой
системой возрастных рейтингов. Нативная поддержка Unity WebGL отсутствует.
Коммерческая монетизация для разработчиков не предусмотрена.

\textbf{Яндекс Игры} --- российская платформа браузерных HTML5-игр с глубокой
интеграцией в экосистему Яндекса. Предоставляет монетизацию через рекламный
CPM-инвентарь. Публикация требует обязательного встраивания Yandex Games SDK
(авторизация, реклама, лидерборды), прохождения верификации и наличия аккаунта
разработчика Яндекса. Unity WebGL без адаптации к SDK не поддерживается, что
существенно ограничивает независимость разработчика и исключает публикацию сборок без
доработки.

Сравнительный анализ ключевых характеристик рассматриваемых платформ и разрабатываемой
системы приведён в таблице~\ref{tab:competitors}.

\begin{landscape}
\begin{table}[p]
    \centering
    \caption{Сравнительный анализ платформ публикации игр}
    \label{tab:competitors}
    \small
    \begin{tabularx}{23cm}{|p{2.9cm}|X|p{2cm}|p{2cm}|p{2.8cm}|X|p{2.8cm}|}
        \hline
        \textbf{Платформа} &
        \textbf{Поддерживаемые форматы} &
        \textbf{Браузерный запуск} &
        \textbf{Unity WebGL} &
        \textbf{Комиссия / модель доступа} &
        \textbf{Монетизация разработчика} &
        \textbf{Барьер входа / модерация} \\ \hline
        itch.io &
        HTML5, WebGL, загружаемые архивы любых форматов &
        Да (iframe) &
        Да (нативно) &
        0--30\% (выбор разработчика) &
        Свободная цена, пожертвования, продажи &
        Низкий: мгновенная публикация, слабая модерация \\ \hline
        GameJolt &
        HTML5, загружаемые файлы &
        Да &
        Частично &
        0\% &
        Отсутствует &
        Низкий: умеренная автоматическая модерация \\ \hline
        Newgrounds &
        HTML5 (Flash устарел) &
        Да &
        Нет &
        0\% &
        Отсутствует &
        Средний: ручная модерация, возрастная классификация \\ \hline
        Яндекс Игры &
        HTML5 + обязательный Yandex Games SDK &
        Да &
        Нет (требует SDK) &
        30\% рекламного дохода &
        Реклама (CPM) &
        Высокий: верификация, SDK, аккаунт Яндекса \\ \hline
        \textbf{Разрабатываемая платформа} &
        Unity WebGL (\texttt{.wasm}, \texttt{.data}, \texttt{.js}, \texttt{.html}) &
        Да (canvas) &
        Да (нативно) &
        0\%, открытый доступ &
        Открытый доступ без комиссии &
        Низкий: регистрация как разработчик, модерация администратором \\ \hline
    \end{tabularx}
\end{table}
\end{landscape}

Таким образом, ни одна из рассмотренных платформ не сочетает в себе все необходимые
характеристики одновременно: открытый бесплатный доступ без внешних SDK, нативную
поддержку Unity WebGL, ролевое разграничение разработчик/игрок и структурированный
каталог с агрегированной статистикой. Это подтверждает целесообразность разработки
специализированной открытой платформы.

\subsection{Анализ технологий веб-разработки и хранения файлов}

Выбор технологического стека обусловлен требованиями к загрузке и браузерному запуску
Unity WebGL сборок, управлению пользователями с ролевой моделью и агрегации игровой
статистики. Ниже приведено последовательное сравнение альтернатив по каждому уровню
архитектуры с обоснованием принятых решений.

\textbf{Выбор серверного фреймворка.} Для реализации серверной части рассматривались
четыре фреймворка: Flask, Django, FastAPI и Express.js. Их сравнение приведено
в таблице~\ref{tab:frameworks}.

\begin{table}[ht!]
    \centering
    \caption{Сравнение серверных фреймворков}
    \label{tab:frameworks}
    \small
    \begin{tabularx}{\textwidth}{|p{2.3cm}|p{1.8cm}|X|X|p{2.2cm}|}
        \hline
        \textbf{Фреймворк} & \textbf{Язык} & \textbf{Преимущества} & \textbf{Недостатки} & \textbf{Решение} \\ \hline
        Flask &
        Python &
        Микрофреймворк: минимальные накладные расходы, полная свобода выбора компонентов, богатая экосистема расширений (Flask-Login, Flask-Migrate, SQLAlchemy) &
        Отсутствие встроенных ORM и admin-панели; всё настраивается вручную &
        \textbf{Выбран} \\ \hline
        Django &
        Python &
        «Батарейки в комплекте»: встроенные ORM, admin-панель, аутентификация, формы, система миграций &
        Избыточная архитектура для небольших проектов, жёсткая структура MVC &
        Не выбран \\ \hline
        FastAPI &
        Python &
        Асинхронная обработка запросов, автоматическая OpenAPI-документация, строгая типизация через Pydantic &
        Молодая экосистема, требует отдельного фронтенд-приложения (REST API), меньше обучающих материалов &
        Не выбран \\ \hline
        Express.js &
        JavaScript &
        Высокая производительность неблокирующего I/O, единый язык с фронтендом, широкое сообщество &
        Другой язык и экосистема (JavaScript вместо Python), менее удобная работа с файлами &
        Не выбран \\ \hline
    \end{tabularx}
\end{table}

Flask выбран как микрофреймворк, позволяющий гибко формировать архитектуру под
конкретные требования~\cite{flask_docs, grinberg2018flask}. Расширения Flask-Login,
Flask-Migrate и SQLAlchemy~\cite{sqlalchemy_docs} закрывают потребности в аутентификации,
миграциях и ORM без избыточности Django.

\textbf{Выбор системы управления базами данных.} Для хранения данных рассматривались
реляционные и нереляционные СУБД. Их сравнение приведено в таблице~\ref{tab:databases}.

\begin{table}[ht!]
    \centering
    \caption{Сравнение систем управления базами данных}
    \label{tab:databases}
    \small
    \begin{tabularx}{\textwidth}{|p{2.5cm}|p{2cm}|X|X|p{2.2cm}|}
        \hline
        \textbf{СУБД} & \textbf{Тип} & \textbf{Преимущества} & \textbf{Недостатки} & \textbf{Решение} \\ \hline
        PostgreSQL &
        Реляционная &
        ACID-транзакции, сложные запросы и агрегации, поддержка JSON-полей, широкий набор индексов, открытый код &
        Требует настройки отдельного сервера &
        \textbf{Выбрана} \\ \hline
        MySQL &
        Реляционная &
        Широкая распространённость, высокая производительность на чтение, простая настройка &
        Исторически слабая поддержка транзакций и внешних ключей (InnoDB решает частично), меньше аналитических функций &
        Не выбрана \\ \hline
        SQLite &
        Реляционная &
        Не требует сервера, нулевая настройка, удобна для разработки и тестирования &
        Ограниченная конкурентная запись, не пригодна для продакшена при нагрузке &
        Не выбрана \\ \hline
        MongoDB &
        Документная (NoSQL) &
        Гибкая схема без миграций, горизонтальное масштабирование &
        Слабая поддержка JOIN и нормализации, более сложные агрегации для рейтингов &
        Не выбрана \\ \hline
    \end{tabularx}
\end{table}

PostgreSQL выбрана как реляционная СУБД с полной поддержкой транзакций, составных
индексов и агрегатных функций, необходимых для расчёта рейтингов и статистики
запусков~\cite{postgresql_docs}. Работа с базой данных осуществляется через
ORM-библиотеку SQLAlchemy, обеспечивающую декларативное описание моделей и
переносимость между СУБД.

\textbf{Выбор подхода к формированию фронтенда.} Для рендеринга пользовательских
страниц рассматривались три подхода: серверный рендеринг (SSR), одностраничные
приложения (SPA) и гибридный HTML-first. Их сравнение приведено в таблице~\ref{tab:frontend}.

\begin{table}[ht!]
    \centering
    \caption{Сравнение подходов к разработке фронтенда}
    \label{tab:frontend}
    \small
    \begin{tabularx}{\textwidth}{|p{2.5cm}|p{2.3cm}|X|X|p{2.2cm}|}
        \hline
        \textbf{Подход} & \textbf{Технология} & \textbf{Преимущества} & \textbf{Недостатки} & \textbf{Решение} \\ \hline
        SSR-шаблоны &
        Jinja2 (Flask) &
        Встроен в Flask, не требует отдельного frontend-проекта и процесса сборки, SEO-дружественность, низкий порог входа &
        Ограниченная клиентская интерактивность без дополнительного JavaScript &
        \textbf{Выбран} \\ \hline
        SPA &
        React / Vue.js &
        Высокая интерактивность, богатый UX, компонентный подход, многократное использование кода &
        Требует отдельного REST/GraphQL API, разделения репозитория, дополнительного процесса сборки и деплоя &
        Не выбран \\ \hline
        HTML-first &
        HTMX &
        Минимум JavaScript, частичная замена фрагментов страницы без перезагрузки, работает поверх SSR &
        Ограниченный набор UI-паттернов, небольшое сообщество, нестандартные атрибуты HTML &
        Не выбран \\ \hline
        Гибридный SSR &
        Next.js &
        SEO, производительность, поддержка React-компонентов, серверные функции &
        JavaScript/Node.js стек, избыточная сложность для CRUD-платформы &
        Не выбран \\ \hline
    \end{tabularx}
\end{table}

Jinja2 выбран как шаблонизатор, встроенный в Flask~\cite{flask_docs}, обеспечивающий
серверный рендеринг HTML-страниц без подключения отдельного frontend-фреймворка. Такой подход соответствует
масштабу проекта и устраняет необходимость разделения на backend-API и
frontend-приложение с раздельным деплоем.

Ключевые технологические решения итоговой архитектуры:

\begin{itemize}
    \item \textbf{Хранение файлов:} игровые сборки Unity WebGL представляют
          собой набор статических файлов (\texttt{.html}, \texttt{.js},
          \texttt{.wasm}, \texttt{.data}), хранимых в файловой системе сервера
          и отдаваемых напрямую через маршруты Flask.
    \item \textbf{Миграции БД:} Flask-Migrate (Alembic) обеспечивает
          версионирование схемы базы данных и применение изменений без потери
          данных.
    \item \textbf{Аутентификация:} сессионная модель на основе Flask-Login с
          хешированием паролей через \texttt{werkzeug.security}.
\end{itemize}

Запуск Unity WebGL в браузере осуществляется через стандартный загрузчик,
поставляемый в составе сборки. Платформа встраивает его в элемент \texttt{canvas}
страницы, не требуя установки дополнительных плагинов.

\subsection{Формулировка цели и задач работы}

На основе проведённого обзора существующих платформ (см. подраздел 1.1) и
анализа технологического стека (см. подраздел 1.2) сформулирована цель работы.

\textbf{Цель работы} --- разработка веб-платформы для публикации, распространения
и запуска браузерных игр, обеспечивающей полный цикл взаимодействия разработчика
и игрока в рамках единого приложения.

Для достижения цели необходимо решить следующие \textbf{задачи}:
\begin{enumerate}
    \item Провести анализ существующих игровых платформ и обосновать выбор
          технологического стека;
    \item Спроектировать модель данных для сущностей User, Game, Comment и Notification;
    \item Реализовать систему аутентификации и ролевого разграничения доступа
          (гость, пользователь, разработчик, администратор);
    \item Реализовать загрузку, хранение и браузерный запуск сборок Unity WebGL;
    \item Реализовать CRUD-интерфейсы для управления играми, комментариями и
          пользователями;
    \item Обеспечить агрегацию игровой статистики (рейтинг, количество запусков)
          и публичный каталог с фильтрацией;
    \item Выполнить деплой приложения и оформить отчётную документацию.
\end{enumerate}
